\section{Results}\label{sec:results}

\subsection{Completeness and execution}
Of 33 request directories, 27 (81.8\%) contain all expected artifacts and six
(18.2\%) contain source only.  All 27 complete sessions report execution status
\code{completed}; execution status is unavailable for the six incomplete
sessions.  The distinction matters: the study cannot determine whether those
six requests failed during compilation, were interrupted, or were retained
before downstream artifacts were written.

\begin{figure}[t]
\centering
\includegraphics[width=\linewidth]{session-outcomes.pdf}
\caption{Artifact completeness and observed outcomes for all 33 request
directories. Source-only observations remain explicit negative/incomplete
results.}
\label{fig:session-outcomes}
\end{figure}

\begin{table}[t]
\centering
\caption{Empirical totals derived from the normalized datasets.}
\label{tab:empirical}
\begin{tabularx}{\linewidth}{@{}XXX@{}}
\toprule
Measure & Derived value & Source \\
\midrule
Complete sessions & 27 & runtime\_summary.json \\
Incomplete sessions & 6 & runtime\_summary.json \\
Normalized sessions & 33 & sessions.csv \\
Complete-session policy violations & 1188 & sessions.csv \\
Counted ledger events & 7155 & event\_counts.csv \\
Prompt-content traces inspected / present & 27 / 0 & runtime\_summary.json \\
Verified evidence bundles & 27 / 27 & verification-results.json \\
\bottomrule
\end{tabularx}

\end{table}

\Cref{fig:session-outcomes} keeps the six source-only directories in view,
while \cref{tab:empirical} reports aggregate totals only where normalized
values exist; together they prevent complete-case counts from becoming a
33-session success claim.

\subsection{Policy inconsistency}
Each of the 27 complete sessions has \code{security_checks=passed} at the
top level.  Nevertheless, each detailed policy object has
\code{policy_passed=false}, \code{review_required=true}, and 44 violations,
all carrying the reason ``unknown policy rule.''  The total is therefore 1,188
policy violations.  These rules cover ATrust, trust-ledger, MCP, and A2A
checks, among others.

We treat this as an inconsistent outcome, not a successful policy evaluation.
The top-level field may describe a narrower check or stale aggregation, but the
artifacts do not establish which.  The appropriate behavior is review and
diagnosis.  This finding illustrates why evidence schemas should preserve
component status and why consumers must not collapse a report to one favorable
label.

\subsection{Event totals and boundaries}
The 27 complete traces contain 7,155 counted ledger events, exactly 265 per
complete session.  Thus the corpus has low structural variation.  Each complete
trace includes 142 \code{PolicyEvaluated} events, 44
\code{PolicyViolation} events, one \code{ExternalProviderExecutionBlocked}
event, one \code{RuntimeNetworkDenied} event, and one
\code{RuntimeSecretsDenied} event.  It also records declaration events for two
agents, six policies, two passports, governance, threat, public-conformance,
release, ATrust, MCP, and A2A metadata.

The repeated structure supports reproducibility of the pipeline but weakens
claims about behavioral diversity.  It does not show how the system responds
to varied prompts, providers, jurisdictions, policy sets, or attacks.  Nor does
a recorded denial independently prove host-level enforcement.

\subsection{Offline bundle verification}
All 27 complete bundles pass the repository's offline verifier (27/27).  Their
bytecode, trace, security-report, and trace-event-ledger digest fields are syntactically valid
SHA-256 identifiers.  The successful result answers RQ3 narrowly: given each
bundle and its relative artifact paths, the verifier can reproduce the
semantic digests, recompute the ledger digest from trace events, check the
SecurityReport LedgerSummary's stored digest against it, and validate
cross-artifact linkage without network access.  The
verifier does not digest or verify \code{session.argx}; source integrity is
outside this result.

It does not answer whether policy approved the request; the detailed policy
result is negative.  It does not authenticate the producer, prove absence of
side effects, or establish semantic correctness of every event.  Integrity and
approval are orthogonal axes.

\subsection{Passports and prompts}
Every complete session reports two passports, Chile as country and
jurisdiction, and Chilean and European Union residency declarations.  These
values show that sovereign metadata survives to runtime evidence.  They do not
show that data was physically stored in either region.

The bounded \code{injected.content} field is present in 0 of the 27 complete,
valid traces.  The six source-only directories have no trace and are not
assessable.  This presence metric is independent of the publication allowlist:
\code{prompt_text} remains null unless explicitly allowlisted even when content
is present.  In this corpus no qualitative prompt themes, quotations,
injection labels, or attack success rates can be reported.  This is
simultaneously a data-minimization benefit and an evaluation limitation.
